% Generated draft: requires mathematical and source review.
\documentclass[11pt]{article}
\usepackage[margin=25mm]{geometry}
\usepackage{fontspec}
\setmainfont{Latin Modern Roman}
\usepackage{amsmath,amssymb,mathtools,unicode-math}
\setmathfont{Latin Modern Math}
\usepackage{enumitem,needspace,xurl,hyperref}
\setlength{\parindent}{0pt}
\setlength{\parskip}{5pt}
\setlength{\emergencystretch}{3em}
\newcommand{\N}{\mathbb{N}}
\newcommand{\Z}{\mathbb{Z}}
\newcommand{\Q}{\mathbb{Q}}
\newcommand{\R}{\mathbb{R}}
\newcommand{\C}{\mathbb{C}}
\newcommand{\F}{\mathbb{F}}
\newcommand{\A}{\mathbb{A}}
\newcommand{\PP}{\mathbb{P}}
\newcommand{\E}{\mathbb{E}}
\newcommand{\eps}{\varepsilon}
\newcommand{\dd}{\,\mathrm{d}}
\begin{document}
\section*{UnsolvedMath ID: 800407}
\textbf{Source problem code:} catalogue entry
\subsection{Definitions and mathematical statement}
Let \(N\ge 2\) and \(d\ge 2\) be integers. The supplied source uses the notation \(\overline{\mu}_d(\mathbb{P}^N)\), but the supplied record does not define this quantity or specify the base field of the projective space \(\mathbb{P}^N\). Consequently, a self-contained definition of the invariant cannot be recovered from this record.

The source reports the bound
\[
\overline{\mu}_d(\mathbb{P}^N)\le d^{-(N-1)}
\qquad\text{for all integers }N\ge 2\text{ and }d\ge 2,
\]
with references [arxiv0908.3835, Proposition 10] and [MR2975153]. This bound is reproduced as a source-reported result.

\textbf{Question 30.} With \(\overline{\mu}_d(\mathbb{P}^N)\) understood according to its intended definition in the cited source, determine its exact value for every pair of integers \(N\ge 2\) and \(d\ge 2\).

\textbf{Unresolved qualification.} The definition of \(\overline{\mu}_d\), including its underlying objects, hypotheses, and any quantifiers or extremal operations, must be supplied before this question becomes a precise, self-contained mathematical problem. The stated upper bound alone does not specify that definition or assert equality.
\subsection{Short English statement}
For each pair of integers \(N,d\ge 2\), determine the exact value of the source's quantity \(\overline{\mu}_d(\mathbb{P}^N)\). The source reports an upper bound of \(d^{-(N-1)}\), but the supplied record omits the definition of the quantity and the base field, so the formulation requires supplementation and human mathematical review.
\subsection{Sources}
\begin{itemize}
\item Supplied local UnsolvedMath registry record: permanent ID \(800407\), exact source code \texttt{catalogue entry}, title ``Open Problems in Arithmetic Dynamics --- Question 30.'' The registry labels the problem \texttt{open}; this is a supplied status label, not an independently verified literature assessment.
\item Benedetto, DeMarco, Ingram, Jones et al., \emph{Current Trends and Open Problems in Arithmetic Dynamics} (2018), Question 30. Supplied source URL: \url{https://arxiv.org/abs/1806.04980}. The record attributes the problem to these authors and gives proposed year 2018.
\item The source statement cites [arxiv0908.3835, Proposition 10] and [MR2975153] for the reported upper bound. Full bibliographic details for these references are not supplied.
\item Supplied provenance: accessed 2026-10-16; extraction method \texttt{arxiv\_latex}; upstream revision \texttt{d371ee917ac9647fbb06cd6fc529deec6f3dda82}. The access date is supplied provenance; no independent source visit or verification is claimed.
\end{itemize}
\end{document}
